%-------------------------------------------
% Yufeng (Alex) Xia - Cover Letter
% Keysight: O-RAN Digital Twin / AI-ML RAN Optimisation Internship
%-------------------------------------------

\documentclass[letterpaper,11pt]{article}

\usepackage[empty]{fullpage}
\usepackage[usenames,dvipsnames]{color}
\usepackage[hidelinks]{hyperref}
\usepackage{fancyhdr}
\usepackage[english]{babel}
\usepackage{lmodern}
\usepackage{parskip}

\pagestyle{fancy}
\fancyhf{}
\fancyfoot{}
\renewcommand{\headrulewidth}{0pt}
\renewcommand{\footrulewidth}{0pt}

\addtolength{\oddsidemargin}{-0.5in}
\addtolength{\evensidemargin}{-0.5in}
\addtolength{\textwidth}{1in}
\addtolength{\topmargin}{-.5in}
\addtolength{\textheight}{1.0in}

\urlstyle{same}
\raggedright
\setlength{\tabcolsep}{0in}
\pdfgentounicode=1

\begin{document}

%----------HEADING (matches CV)-----------------
\begin{tabular*}{\textwidth}{l@{\extracolsep{\fill}}r}
  \textbf{\href{https://www.linkedin.com/in/yufeng-xia-12648a25b}{\Large Yufeng (Alex) Xia}} & Email: \href{mailto:yufeng.xia@ed.ac.uk}{yufeng.xia@ed.ac.uk} \\
  \href{https://www.linkedin.com/in/yufeng-xia-12648a25b}{linkedin.com/in/yufeng-xia-12648a25b} & Mobile: \href{tel:+447918155093}{(+44) 7918155093} \\
  Edinburgh, UK & \href{https://xyf2002.github.io/}{xyf2002.github.io} \\
\end{tabular*}

\vspace{6pt}
\hrule
\vspace{14pt}


\vspace{8pt}

Keysight Technologies \\
South Queensferry, Edinburgh, United Kingdom

\vspace{10pt}

Dear Hiring Manager,

\vspace{4pt}

I am writing to apply for the internship on O-RAN digital twins and AI/ML-driven
RAN optimisation. My MSc by Research at the University of Edinburgh, supervised
by Prof.\ Mahesh K.\ Marina, sits on precisely the two halves that this role
combines. I build the digital twin and emulation platforms on which RAN
algorithms are developed and validated, and I design the near-real-time policies
that decide how those platforms use shared compute. In my work these have never
been two separate tasks, which is why a role that asks for both is the one I want.

My thesis project, Flux-DT, is an elastic digital twin runtime for AI-RAN.
Instead of observing a network from the outside, the twin runs inside the radio
path of a live OpenAirInterface 5G stack, so that real signalling and user
traffic flow through a ray-traced channel model computed across GPU workers. On
top of it I designed a control loop that reacts to radio-driven resource pressure
on the timescale the radio imposes: it relocates work to less contended machines,
lowers simulation fidelity only within bounds that live radio procedures
tolerate, and suspends and resumes the twin when no further reduction is safe.
The placement decision itself is a joint objective across co-channel cells that I
solve with a greedy heuristic, because the decision has to be made faster than an
exact solver could return.
Getting this right meant treating the live RAN as the priority workload at every
step, and treating the twin's own state as something that must survive being
interrupted. I understand these to be the same constraints that a production
digital twin platform has to respect.

I have also validated systems of this kind at scale. On EmuRAN, a cloud-based
RAN emulation system conditionally accepted at ACM MobiCom 2026, I owned the
evaluation: I deployed and operated the emulator on public cloud at 130 instances
and 6,400 CPU cores, emulating 10,000 mobile devices and 10,000 base stations,
and I designed the fidelity and scalability experiments that established it
behaves like a real RAN rather than merely running at scale. My machine learning
systems work is closely connected to this. Morphling, an emulator for evaluating
distributed machine learning training on heterogeneous edge devices, received the
Best Paper Award at EdgeSys 2026, and on Weaver I evaluated a controller that
decides how much accelerator capacity an AI workload may harvest from shared RAN
infrastructure without degrading radio performance. Alongside this, my C and C++
experience comes from compiler research, where I built an LLVM-based
auto-quantization framework and evaluated it on embedded ARM devices.

What draws me to Keysight specifically is that this role puts algorithm
development and platform work in the same place. Developing AI and machine
learning algorithms for prioritised O-RAN use cases, and then validating them on
digital twins and emulation platforms instead of in a conventional lab, is
exactly the loop I have been building the tools for, and I would much rather
close that loop against real equipment than in simulation alone. The near-real-time
use cases interest me most --- traffic steering, load balancing and energy saving,
where a decision has to respect interference and load conditions within a deadline
rather than converge eventually. My control-loop work has been at exactly that
timescale, though outside a RIC; the natural next step is to develop such logic as
an xApp against the near-RT RIC and validate it on a digital twin, which is what
this role does. Applying AI and machine learning to the optimisation of cellular
network resources, in the context of 5G and future 6G radio access, is the
direction I intend to build my career in.

I graduate in September 2026 and am available for a full-time internship from
then. I am already based in Edinburgh, so I can work on site without relocating.
I would welcome the opportunity to discuss how I could contribute to the team,
and I am happy to talk through any of the systems above in more detail.

\vspace{10pt}

Yours faithfully,

\vspace{14pt}

\textbf{Yufeng Xia}

\end{document}
